% ============================================================
% PARTE IV - PROGRESSIONI DI TRIADI
% CAPITOLO 15 - LE PROGRESSIONI FONDAMENTALI
% ============================================================
\part{Progressioni di triadi}

\chapter{Le progressioni fondamentali}
\label{chap:progressioni_armoniche}

Funzioni, cadenze e Regola dell'Ottava permettono di armonizzare brevi tratti di basso. Per i tratti lunghi o ripetitivi serve un altro strumento: la \textbf{progressione}, una successione di accordi che si ripete trasposta.

\section{Che cos'è una progressione}
\label{sec:definizione_progressione}

\begin{definizione}[Progressione armonica]
	\label{def:progressione_armonica}
	Una \textbf{progressione armonica} (o \textbf{sequenza armonica}) è una successione di accordi ripetuta trasportata su gradi diversi della scala. La ripetizione crea un movimento armonico coerente e prevedibile, con cui si costruiscono frasi lunghe o si modula.
\end{definizione}

Una progressione si descrive con quattro elementi:
\begin{itemize}
	\item il \textbf{modello}, una breve successione di accordi (in genere 2-4);
	\item la \textbf{trasposizione}, cioè la ripetizione del modello su gradi diversi;
	\item la \textbf{direzione}, verso l'acuto (ascendente) o verso il grave (discendente);
	\item il \textbf{numero di ripetizioni}, in genere 2-4, per non cadere nella monotonia.
\end{itemize}

Le progressioni danno una soluzione sistematica ai passaggi scalari lunghi, creano movimento e facilitano la modulazione ai toni vicini. Perché l'orecchio le segua con tanta facilità è detto nell'Approfondimento \ref{phys:perche_progressioni}; le progressioni si possono anche leggere come percorsi tra livelli di tensione (Approfondimento \ref{phys:energia_armonica}).

\begin{errore}[Progressioni troppo lunghe]
	\label{err:progressioni_lunghe}
	Applicare il modello per troppe ripetizioni, per esempio una progressione per quinte discendenti che occupa tre misure intere, rende la musica meccanica e prevedibile. Conviene interrompere la progressione dopo 2-3 cicli e variare: chiudere la frase con una cadenza, cambiare tipo di progressione oppure inserire un accordo di passaggio o una nota pedale. Raramente si superano le 3-4 ripetizioni consecutive dello stesso modello.
\end{errore}

\section{Progressione per quinte discendenti}
\label{subsec:progressione_quinte}

La progressione per quinte discendenti è una delle più naturali del sistema tonale, perché replica il movimento cadenzale V-I su ogni grado della scala. Il basso si muove per quinte discendenti, o quarte ascendenti, e la numerica è 5 su tutti gli accordi.

\begin{regola}[Progressione per quinte discendenti]
	\label{reg:progressione_quinte_disc}
	Ogni accordo è seguito da un accordo la cui fondamentale sta una quinta sotto (o una quarta sopra). Lo schema è:

	\textbf{modo maggiore:} I - IV - VII - III - VI - II - V - I

	\textbf{modo minore:} i - iv - VII - III - VI - ii$^\circ$ - V - i

	Caratteristiche:
	\begin{itemize}
		\item il basso procede per \textbf{quarte ascendenti}, equivalente intervallare della quinta discendente;
		\item gli accordi sono in \textbf{stato fondamentale} (numerica 5, o nessuna cifra);
		\item per mantenere fluido il movimento del basso si può inserire occasionalmente un rivolto, per esempio il \sigla{VII}{6}{};
		\item la sensazione è di continua risoluzione verso la tonica.
	\end{itemize}
\end{regola}

Nella condotta delle voci tra due accordi in stato fondamentale con basso che salta di quinta, le tre voci superiori si muovono in \textbf{moto contrario} rispetto al basso, oppure conservano una o due note comuni, in modo da evitare quinte e ottave parallele.

\begin{esempio}[Progressione per quinte discendenti in Do maggiore]
	\label{ex:progressione_quinte_disc}
	Alternanza di stato fondamentale (5) e primo rivolto (6) per mantenere il basso in movimento per grado congiunto o per piccoli salti.
	\musica{es_cap17_01}
\end{esempio}

\begin{esempio}[Progressione per quinte discendenti (tutti 5) in Sol maggiore]
	\label{ex:progressione_5_sol}
	Il basso scende per quinte (Sol-Do-Fa$\sharp$-Si-Mi-La-Re-Sol). Per non appesantire la condotta si usa qualche rivolto, come il \sigla{VII}{6}{}.
	\musica{es_cap17_11}
\end{esempio}

L'ordine dei gradi nella progressione ha una descrizione matematica (Approfondimenti \ref{math:cerchio_quinte} e \ref{math:gruppo_z7}, Appendice \ref{app:approfondimenti}).

\section{Progressione per terze discendenti}
\label{subsec:progressione_terze_disc}

La progressione per terze discendenti ha un movimento più morbido rispetto a quella per quinte.

\begin{regola}[Progressione per terze discendenti]
	\label{reg:progressione_terze_disc}
	Ogni accordo scende di terza (maggiore o minore) rispetto al precedente. Nel modo maggiore la forma comune è
	I - VI - IV - II - VII - V - I.

	Caratteristiche:
	\begin{itemize}
		\item il movimento è più graduale e «dolce» di quello per quinte;
		\item si usa spesso in passaggi espressivi o malinconici;
		\item il basso si muove per terze o, più spesso, per grado congiunto con rivolti appropriati.
	\end{itemize}
\end{regola}

\begin{esempio}[Progressione per terze discendenti in Do maggiore]
	\label{ex:progressione_terze_disc}
	La sequenza scende attraverso la scala, con una sensazione di dolce malinconia.
	\musica{es_cap17_02}
\end{esempio}
